% chapters/introduction.tex — Introduction chapter (narrative text only)
\chapter{Introduction}
\label{chap:introduction}

The emergence of large language models (\LLM{}s) capable of autonomous,
multi-step reasoning has opened a new research frontier: \emph{agentic
systems} that perceive their environment, plan sequences of actions, and
interact with external tools to accomplish complex goals~\cite{wei2026agenticreasoning}.
Unlike single-turn question answering, agentic \LLM{}s must maintain context
across many steps, recover from errors, and coordinate with other agents or
subsystems~\cite{ke2025reasoningfrontiers}.
\index{agent}
\index{large language model}

\section{Motivation}
\label{sec:intro-motivation}

Classical approaches to agent design relied on hand-crafted planners and
symbolic state representations.  Modern \LLM{}-based agents replace the
symbolic backbone with a pre-trained model whose internal representations
encode broad world knowledge~\cite{wei2026agenticreasoning}.
This shift enables rapid adaptation to novel tasks but introduces new
challenges around reliability, memory, and grounding.

\section{Scope and Contributions}
\label{sec:intro-scope}

This survey covers four core capabilities of reasoning-centric agentic
\LLM{}s: planning and task decomposition, memory and state persistence,
retrieval-augmented grounding, tool use, and multimodal reasoning.
We also discuss evaluation frameworks and open research directions.

\section{Overview in Hebrew / סקירה בעברית}
\label{sec:intro-bidi}

This section demonstrates correct bidirectional text rendering under
LuaLaTeX with \texttt{polyglossia} and the David~CLM Hebrew font.
\index{bidirectional text}
\index{סוכן@\texthebrew{סוכן} (agent)}

\medskip
\begin{hebrew}
מערכות \textenglish{LLM} אגנטיות מאפשרות ביצוע משימות מורכבות
על ידי תכנון, זיכרון ושימוש בכלים חיצוניים.
הסקר הזה סוקר את הגישות המובילות בתחום:
\textenglish{planning}, \textenglish{memory}, \textenglish{retrieval},
שימוש בכלים (\textenglish{tool use}), וחישוב רב-מודאלי
(\textenglish{multimodal reasoning}).
מחקרים כמו \textenglish{\cite{wei2026agenticreasoning}}
מניחים את היסודות לפיתוח סוכנים חכמים יותר.
\end{hebrew}
\medskip

The paragraph above is typeset right-to-left, with embedded English
technical terms switching back to left-to-right rendering automatically
via \texttt{polyglossia}.

\section{Chapter Outline}
\label{sec:intro-outline}

Chapter~\ref{chap:planning} surveys planning and task decomposition.
Chapter~\ref{chap:memory} covers memory architectures.
Chapter~\ref{chap:retrieval} discusses retrieval-augmented reasoning.
Chapter~\ref{chap:tooluse} examines tool-use policies.
Chapter~\ref{chap:multimodal} addresses multimodal agentic systems.
Chapter~\ref{chap:evaluation} reviews evaluation benchmarks.
Chapter~\ref{chap:conclusion} summarises findings and future directions.
